% Source printed pp. 81-83.
\section{Grothendieck 拓扑与景}\label{src:3.2}
\begin{epigraph}
引入数字 $0$ 或群的概念也曾是一般的抽象废话, 而数学几千年来多少一直停滞不前, 就是因为没有人迈出这样幼稚的一步...
\epigraphauthor{Alexander Grothendieck, 致 Ronald Brown 的信}
\end{epigraph}
为利用上一节的下降性质, 我们将拓扑空间的概念加以推广, 从而将平展, fppf 或 fpqc 态射看作开集, 并在这些广义拓扑中表述层与叠的公理. Grothendieck 拓扑是在~\cite[定义 II.1.3]{SGA4}~中引入的. 我们的阐述遵循~\cite{Art62}, \cite{Vis05}, \cite[§2]{Ols16}~及~\cite[Tag 00UZ]{SP}.

\subsection{定义与最初的例子}\label{src:3.2.1-subsection}
\begin{definition}[景]\label{src:3.2.1}
范畴 $\mathcal S$ 上的\newterm{Grothendieck 拓扑}{Grothendieck topology}由以下数据组成: 对每个对象 $X\in\mathcal S$, 有一个集合 $\Cov(X)$, 其元素为 $X$ 的覆盖, 即 $\mathcal S$ 中的态射族 $\{X_i\to X\}_{i\in I}$. 要求:
\begin{enumerate}[(1)]
\item (恒等) 若 $X'\to X$ 是同构, 则 $(X'\to X)\in\Cov(X)$.
\item (限制) 若 $\{X_i\to X\}_{i\in I}\in\Cov(X)$ 且 $Y\to X$ 是态射, 则纤维积 $X_i\times_X Y$ 在 $\mathcal S$ 中存在, 且态射族 $\{X_i\times_X Y\to Y\}_{i\in I}$ 属于 $\Cov(Y)$.
\item (复合) 若 $\{X_i\to X\}_{i\in I}\in\Cov(X)$, 且对每个 $i\in I$ 都有 $\{X_{ij}\to X_i\}_{j\in J_i}\in\Cov(X_i)$, 则 $\{X_{ij}\to X_i\to X\}_{i\in I,j\in J_i}\in\Cov(X)$.
\end{enumerate}
\newterm{景}{site}是带有 Grothendieck 拓扑的范畴 $\mathcal S$.
\end{definition}
\begin{caution}\label{src:3.2.2}
定义要求 $\Cov(X)$ 是集合, 但在许多重要情形中, 例如大平展景, 这一点并不成立. 本书将忽略集合论问题. 通常可以通过选取适当的子范畴来处理这些问题, 这个子范畴应包含所有需要考虑的态射, 且定义相同的层范畴; 这一般容易做到, 但颇为繁琐. 见~\cite[Tag 00VI]{SP}.
\end{caution}
\begin{example}[拓扑空间]\label{src:3.2.3}
若 $X$ 是拓扑空间, 记 $\Op(X)$ 为开集 $U\subseteq X$ 的范畴, 其中存在唯一的态射 $U\to V$ 当且仅当 $U\subseteq V$. 我们称开子集族 $\{U_i\}_{i\in I}$ 为 $U$ 的覆盖 (即 $\Cov(U)$ 中的元素), 若 $U=\bigcup_{i\in I}U_i$. 这定义了 $\Op(X)$ 上的 Grothendieck 拓扑. 特别地, 若 $X$ 是概形, 则 $X$ 上的 Zariski 拓扑定义一个景 $X_{\mathrm{Zar}}$, 称为 $X$ 上的\newterm{小 Zariski 景}{small Zariski site}.
\end{example}
将 Zariski 开浸入替换为平展态射, 就得到小平展景.
\begin{example}[小平展景]\label{src:3.2.4}
若 $X$ 是概形, 则 $X$ 上的\newterm{小平展景}{small étale site}是由平展态射 $U\to X$ 组成的范畴 $X_{\text{ét}}$, 其中态射 $(U\to X)\to(V\to X)$ 就是 $X$-态射 $U\to V$ (它必然是平展的). 换言之, $X_{\text{ét}}$ 是 $\Sch/X$ 中由 $X$ 上的平展概形组成的全子范畴. 对象 $(U\to X)\in X_{\text{ét}}$ 的覆盖是平展态射族 $\{U_i\to U\}$, 使得 $\coprod_i U_i\to U$ 满射. 后面将介绍代数空间或 Deligne--Mumford 叠的小平展景 $X_{\text{ét}}$ (定义~\bookref{5.1.1}), 用以定义 $X$ 上的层.
\end{example}

\subsection{大景}\label{src:3.2.2-subsection}
\begin{example}[大平展景]\label{src:3.2.5}
\newterm{大平展景}{big étale site}是范畴 $\Sch$, 其中概形 $U$ 的覆盖是 $\Sch$ 中的平展态射族 $\{U_i\to U\}$, 使得 $\coprod_i U_i\to U$ 满射. 我们将这个景记为 $\Sch_{\text{ét}}$.
\end{example}
大平展景 $\Sch_{\text{ét}}$ 是本书最重要的景, 用来定义本书最核心的概念: 代数空间是 $\Sch_{\text{ét}}$ 上在平展局部为概形的层 (定义~\bookref{4.1.2}), 而代数叠是 $\Sch_{\text{ét}}$ 上在光滑局部为概形的叠 (定义~\bookref{4.1.6}). 还有若干类似的景.
\begin{example}[大拓扑景]\label{src:3.2.6}
设 $\mathsf{Top}$ 是拓扑空间范畴. $U\in\mathsf{Top}$ 的覆盖是开子空间族 $\{U_i\hookrightarrow U\}_{i\in I}$, 使得 $U=\bigcup_{i\in I}U_i$.
\end{example}
\begin{example}[大 Zariski 景]\label{src:3.2.7}
将例~\bookref{3.2.5}~中的平展态射替换为开浸入, 就定义了大 Zariski 景 $\Sch_{\mathrm{Zar}}$.
\end{example}
\begin{example}[大 fppf 景]\label{src:3.2.8}
概形态射是 fppf 的, 是指它满射, 平坦且局部有限表现 (见定义~\bookref{A.2.20}). 大 fppf 景 $\Sch_{\mathrm{fppf}}$ 是概形范畴 $\Sch$, 其中覆盖 $\{U_i\to U\}$ 是满足 $\coprod_i U_i\to U$ 为 fppf 态射的态射族.
\end{example}
大 fppf 景 $\Sch_{\mathrm{fppf}}$ 在模理论中也很重要. 在~\cite{SP}~中, 代数空间与代数叠使用 $\Sch_{\mathrm{fppf}}$ 而非 $\Sch_{\text{ét}}$ 来定义, 但这两个景仍然给出等价的概念~\cite[Tag 076U]{SP}. 本书在 \S\bookref{7.3} 中使用大 fppf 景讨论 gerbe 以及非光滑群概形的商叠.
\begin{example}[大 fpqc 景]\label{src:3.2.9}
概形态射是 fpqc 的, 是指它满射, 平坦, 且靶端的每个拟紧子集都是某个拟紧子集的像 (见定义~\bookref{A.2.20}). 大 fpqc 景 $\Sch_{\mathrm{fppf}}$ 是范畴 $\Sch$, 其中覆盖 $\{U_i\to U\}$ 是满足 $\coprod_i U_i\to U$ 为 fpqc 态射的态射族.
\end{example}
\begin{caution}\label{src:3.2.10}
定义 fpqc 景存在严重的集合论问题, 原因是 fpqc 覆盖太多: 对任意非零环 $R$, 不存在一个由 $\Spec R$ 的 fpqc 覆盖组成的集合, 能加细每个 fpqc 覆盖; 见~\cite[Tags 0BBK, 03NV]{SP}. 若忽略集合论问题来定义大 fpqc 景 (正如我们刚才所做的!), 则有些预层不存在层化~\cite[定理 $5.5$]{Wat75}. 幸运的是, 本书并不需要 fpqc 景. 另一方面, fpqc 拓扑中的层或叠这一概念是良定义的, 因而可以用来表述一般命题 (不过我们通常只使用平展情形).
\end{caution}
\begin{example}[光滑平展景]\label{src:3.2.11}
概形 $X$ 上的\newterm{光滑平展景}{lisse-étale site} $X_{\text{lis-ét}}$ 是由 $X$ 上的光滑概形组成的范畴, 其中 $X_{\text{lis-ét}}$ 中的态射是 $X$ 上的概形态射 (不一定光滑). $X$-概形 $U$ 的覆盖 $\{U_i\to U\}$ 是 $X$-态射族, 使得 $\coprod_i U_i\to U$ 满射且平展.
\end{example}
我们将在定义~\bookref{7.1.1}~中引入代数叠的光滑平展景, 用以定义拟凝聚层.
\begin{example}[限制范畴与限制景]\label{src:3.2.12}
若 $\mathcal S$ 是范畴且 $S\in\mathcal S$, 定义\newterm{限制范畴}{restricted category} (有时也称为局部化范畴) 为范畴 $\mathcal S/S$, 其对象是 $\mathcal S$ 中的映射 $T\to S$. 态射 $(T'\to S)\to(T\to S)$ 是 $S$ 上的映射 $T'\to T$. 若 $\mathcal S$ 是景, 则 $\mathcal S/S$ 也成为景, 其中 $\mathcal S/S$ 中 $T\to S$ 的覆盖就是 $\mathcal S$ 中 $T$ 的覆盖 $\{T_i\to T\}$. 将这一构造应用于概形 $S$, 得到大 Zariski, 平展, fppf 与 fpqc 景的相对版本 $(\Sch/S)_{\mathrm{Zar}}$, $(\Sch/S)_{\text{ét}}$, $(\Sch/S)_{\mathrm{fppf}}$ 及 $(\Sch/S)_{\mathrm{fpqc}}$.
\end{example}
\begin{example}[仿射概形范畴上的 Grothendieck 拓扑]\label{src:3.2.13}
文献中有时会使用仿射概形范畴上的大景 $\mathsf{AffSch}_{\mathrm{Zar}}$, $\mathsf{AffSch}_{\text{ét}}$ 与 $\mathsf{AffSch}_{\mathrm{fppf}}$. 它们定义的层范畴与 $\Sch$ 上相应大景的层范畴相同.
\end{example}
